\chapter{Investigation of Interface Favourability} \label{chapter:computational_investigation}

% ---------- %

A key step in training a Machine-Learned Interface (MLI) model is building a dataset of physically plausible and
energetically reasonable heterostructures.

Ideally, such a dataset would be generated entirely using density functional theory (DFT), but the cost of relaxing and
evaluating thousands of configurations makes this impractical.

Instead, this chapter investigates whether a machine-learned potential (MLP), specifically the MACE model, can assist in
rapidly generating and assessing interface structures before applying DFT.

% ---------- %

The focus here is not to validate MLPs as a replacement for DFT, but to test their usefulness as part of a larger
data-generation pipeline.

If MACE is able to help identify realistic and favourable structures, it can reduce the number of expensive DFT
calculations needed and enable faster construction of training data for the MLI model.

% ---------- %

The chapter is structured as follows: Section 1 describes the generation of the dataset, the chemical systems
considered, and the use of ARTEMIS for constructing candidate interfaces.

Section 2 analyses how well the MACE model reproduces trends in DFT-calculated interface energetics.

Section 3 discusses and summarises analysis and findings.

% ---------- %
